\chapter{Materials and Methods}
\label{ch:materials_methods}

\section{Collection and Preparation of the Food Waste}
\label{sec:col_prep_FW}

\Glsentrylongauto{fw} was collected from the canteen at the \glsauto{ipb}, prioritizing residues with high proportions of rice, pasta, and bread to maximize the sugar content and minimize the nitrogen content. It was collected \qty[round-mode = places, round-precision = 0]{3235.04}{\gram} of waste then it was diluted with approximately \qty{2.5}{\litre} of distilled water and homogenised using a blender to obtain a \glsaut{fw} suspension with \qty[round-mode = places, round-precision = 0]{5525.55}{\gram} total. It was then stored in a refrigerator at \qty{4}{\celsius} to prevent spoilage and suppress microbial and fungal activity \parencite{degueurce2020}.

\section{Crude Glycerol Collection}
\label{sec:gly_col}

A \qty{6}{\litre} bottle of \glsauto{cgl}, obtained from previous biodiesel production experiments at \glsauto{ipb}, was provided for this study. \glsauto{cgl} from biodiesel typically contains impurities such as catalyst residues \parencite{attarbachi2023}, which may influence its behaviour in anaerobic digestion.

\section{Inoculum Collection}
\label{sec:ino_col}

The inoculum was collected from an anaerobic digester from Águas do Norte wastewater treatment plant in Bragança. This ensured the availability of a microbial population for the anaerobic digestion \parencite{posmanik2020}.

\section{Physical and Chemical Characterisation}
\label{sec:phy_chem_char}

The \glsauto{ts}, \glsauto{vs}, \glsauto{fs}, and density were determined for the \glsauto{fw}, \glsauto{cgl}, and inoculum. \Glsentrylongauto{toc} and \glsauto{tn} were determined for the \glsauto{fw} and \glsauto{cgl}.

\subsection{TS and VS Determination}
\label{subsec:ts_vs_determination}

\Glsentrylongauto{ts} and \glsentrylongauto{vs} were measured using a sequential gravimetric method based on Standard Methods 2540 B and 2540 E \parencite{apha2023}. The method involves controlled drying and calcination steps to quantify \glsauto{ts} and \glsauto{vs} fractions.

Initially, porcelain crucibles were placed in a muffle furnace at \qty{550}{\celsius} for \qty{1}{\hour}. For cooling, they were transferred to an oven at \qty{105}{\celsius} for \qty{15}{\minute} and then moved to a desiccator until they reached room temperature; their empty masses were recorded. For each material, a sample was placed into three crucibles and the wet mass was measured. The crucibles were then dried in an oven at \qty{105}{\celsius} for \qty{1}{\hour}, the process was repeated until constant mass. After drying, the crucibles were cooled in a desiccator to room temperature and weighed again to determine \glsauto{ts}.

The same crucibles were subsequently placed in a muffle furnace at \qty{550}{\celsius} for \qty{50}{\minute}. Afterwards they were transferred to an oven at \qty{105}{\celsius} for \qty{15}{\minute} and then moved to a desiccator until they reached room temperature; the final masses were recorded to calculate \glsauto{vs}.

\Glsentrylongauto{ts}, \Glsentrylongauto{vs} and \Glsentrylongauto{fs} were calculated as follows:
\begin{equation}
    \mathrm{TS} = \frac{M_{105} - M_{\text{crucible}}}{V_{\text{sample}}} \quad \left[\unit{\gram\per\litre}\right]
    \label{eq:TS}
\end{equation}
\begin{equation}
    \mathrm{VS} = \frac{M_{105} - M_{550}}{V_{\text{sample}}} \quad \left[\unit{\gram\per\litre}\right]
    \label{eq:VS}
\end{equation}
\begin{equation}
    \mathrm{FS} = \frac{M_{550} - M_{\text{crucible}}}{V_{\text{sample}}} \quad \left[\unit{\gram\per\litre}\right]
    \label{eq:FS}
\end{equation}
where $M_{\text{wet}}$ is the mass of crucible plus wet sample, $M_{105}$ is the mass after drying at \qty{105}{\celsius}, $M_{550}$ is the mass after ignition at \qty{550}{\celsius}, $M_{\text{crucible}}$ is the mass of the empty crucible and $V_{\text{sample}}$ is the volume of the sample. The volume is in \unit{\litre} and the mass is in \unit{\gram} (\Cref{tab:solids}).

\begingroup
\tablesetup
\footnotesize
\setlength{\tabcolsep}{3pt}
\begin{landscape}
    \begin{ThreePartTable}
        \begin{xltabular}{\linewidth}{@{} C C C C C C C C c C C c @{}}

        \caption{\Glsentrylongauto{ts}, \Glsentrylongauto{vs} and \Glsentrylongauto{fs} values}
        \label{tab:solids} \\

        \toprule
        \multicolumn{2}{c}{\thead{Sample}} &
        \thead{$M_{\text{crucible}}$\\(\unit{\gram})} &
        \thead{$M_{\text{wet}}$\\(\unit{\gram})} &
        \thead{$M_{\text{sample}}$\\(\unit{\gram})} &
        \thead{$V_{\text{sample}}$\\(\unit{\litre})} &
        \thead{$M_{105}$\\(\unit{\gram})} &
        \thead{\glsauto{ts}\\(\unit{\gram\per\litre})} &
        \thead{Average \glsauto{ts}\\(\unit{\gram\per\litre})} &
        \thead{$M_{550}$\\(\unit{\gram})} &
        \thead{\glsauto{vs}\\(\unit{\gram\per\litre})} &
        \thead{Average \glsauto{vs}\\(\unit{\gram\per\litre})} \\
        \midrule
        \endfirsthead

        \multicolumn{12}{l}{\textit{\Cref{tab:solids} continued}} \\
        \toprule
        \multicolumn{2}{c}{\thead{Sample}} &
        \thead{$M_{\text{crucible}}$\\(\unit{\gram})} &
        \thead{$M_{\text{wet}}$\\(\unit{\gram})} &
        \thead{$M_{\text{sample}}$\\(\unit{\gram})} &
        \thead{$V_{\text{sample}}$\\(\unit{\litre})} &
        \thead{$M_{105}$\\(\unit{\gram})} &
        \thead{\glsauto{ts}\\(\unit{\gram\per\litre})} &
        \thead{Average \glsauto{ts}\\(\unit{\gram\per\litre})} &
        \thead{$M_{550}$\\(\unit{\gram})} &
        \thead{\glsauto{vs}\\(\unit{\gram\per\litre})} &
        \thead{Average \glsauto{vs}\\(\unit{\gram\per\litre})} \\
        \midrule
        \endhead

        \bottomrule
        \endlastfoot

        \directlua{emit_grouped_table("data/processed/solids.csv", 12)}

        \end{xltabular}
    \end{ThreePartTable}
\end{landscape}
\endgroup

\subsection{Density Determination}
\label{subsec:d_determination}

Density was determined by a gravimetric volumetric method using graduated cylinders. Dry cylinders were weighed empty, then filled with distilled water at \qty{20}{\celsius} and weighed again to verify the reading accuracy (density of water at $\rho_{20\,°C} = \qty{998.21}{\gram\per\litre}$). After drying, each cylinder was filled with a known volume of sample and weighed. The density of each material was calculated using
\begin{equation}
    \rho = \frac{M_{\text{total}} - M_{\text{cylinder}}}{V_{\text{sample}}} \quad \left[\unit{\gram\per\litre}\right]
    \label{eq:density}
\end{equation}
where $M_{\text{total}}$ is the mass of cylinder plus sample, $M_{\text{cylinder}}$ is the empty cylinder mass, and $V_{\text{sample}}$ is the measured volume, the masses were measured in \unit{\gram} and the volume in \unit{\litre} (\Cref{tab:density}).

\begingroup
\tablesetup
\footnotesize
\setlength{\tabcolsep}{3pt}
\begin{ThreePartTable}
    \begin{xltabular}{\linewidth}{@{} C C C C C C c @{}}
        \caption{Results of desnsity for substrates and inoculum}
        \label{tab:density} \\

        \toprule
        \multicolumn{2}{c}{\thead{Sample}} &
        \thead{$M_{\text{cylinder}}$ (\unit{\gram})} &
        \thead{$M_{\text{total}}$ (\unit{\gram})} &
        \thead{$V_{\text{sample}}$ (\unit{\litre})} &
        \thead{Density\\(\unit{\gram\per\litre})} &
        \thead{Average Density\\(\unit{\gram\per\litre})} \\
        \midrule
        \endfirsthead

        \multicolumn{7}{l}{\textit{\Cref{tab:density} continued}} \\
        \toprule
        \multicolumn{2}{c}{\thead{Sample}} &
        \thead{$M_{\text{cylinder}}$ (\unit{\gram})} &
        \thead{$M_{\text{total}}$ (\unit{\gram})} &
        \thead{$V_{\text{sample}}$ (\unit{\litre})} &
        \thead{Density\\(\unit{\gram\per\litre})} &
        \thead{Average Density\\(\unit{\gram\per\litre})} \\
        \midrule
        \endhead

        \bottomrule
        \endlastfoot

        \directlua{emit_grouped_table("data/processed/density.csv", 7)}

    \end{xltabular}
\end{ThreePartTable}
\endgroup

\subsection{Determination of TOC and TN}
\label{subsec:toc_tn_determination}

\Glsauto{toc}, \glsauto{tc} and \glsauto{tn} were measured using a Shimadzu TOC-5000A analyser. Prior to analysis, the \glsauto{fw} was diluted \qty{400}{\dil} and the \glsauto{cgl} was diluted \qty{800}{\dil} to bring the concentrations within the linear range of the instrument. The results are on the \Cref{tab:toc}.

\begingroup
\tablesetup
\footnotesize
\setlength{\tabcolsep}{3pt}
\begin{ThreePartTable}
    \begin{xltabular}{\linewidth}{@{} C C C C C C C @{}}

        \caption{\Glsentrylongauto{tc} and \Glsentrylongauto{tn} for substrates}
        \label{tab:toc} \\

        \toprule
        \multicolumn{2}{c}{\thead{Sample}} &
        \thead{\GLS{toc}\\(\unit{\milli\gram\per\litre})} &
        \thead{\GLS{tc}\\(\unit{\milli\gram\per\litre})} &
        \thead{Average \GLS{tc}\\(\unit{\milli\gram\per\litre})} &
        \thead{\GLS{tn}\\(\unit{\milli\gram\per\litre})} &
        \thead{Average \GLS{tn}\\(\unit{\milli\gram\per\litre})} \\
        \midrule
        \endfirsthead

        \multicolumn{7}{l}{\textit{\Cref{tab:toc} continued}} \\
        \toprule
        \multicolumn{2}{c}{\thead{Sample}} &
        \thead{\GLS{toc}\\(\unit{\milli\gram\per\litre})} &
        \thead{\GLS{tc}\\(\unit{\milli\gram\per\litre})} &
        \thead{Average \GLS{tc}\\(\unit{\milli\gram\per\litre})} &
        \thead{\GLS{tn}\\(\unit{\milli\gram\per\litre})} &
        \thead{Average \GLS{tn}\\(\unit{\milli\gram\per\litre})} \\
        \midrule
        \endhead

        \bottomrule
        \endlastfoot

        \directlua{emit_grouped_table("data/processed/toc.csv", 7)}
    \end{xltabular}
\end{ThreePartTable}
\endgroup

\section{Experiment Setup}
\label{sec:experiment_configuration}

The reactor had a working volume of \qty{1}{\litre} and was operated in continuous mode, under anaerobic conditions. Mixing was provided by a magnetic stir bar, and the overflow outlet was kept closed with a valve to maintain an anaerobic environment. This experimental setup is shown in \Cref{graph:reactor}:

\begin{figure}[H]
    \centering
    \resizebox{\textwidth}{!}{%
        \includestandalone[mode=tex]{assets/graphs/reactor-model}
    }
    \caption{Schematic of the experimental setup used}
    \label{graph:reactor}
\end{figure}

The reactor was kept at \qty{35}{\celsius}, close to the optimal temperature for mesophilic microorganisms \parencite{choorit2007}. The \glsauto{olr} was set with a max of \qty{4}{\kilo\gram\VS\per\cubic\metre\per\day}, with a \glsauto{hrt} of \num{30} days, conditions for stable co-digestion of \glsauto{fw} and \glsauto{gl} \parencite{tomczak2025}. The pH was maintained around \num{7} by adding \glsauto{khco3} into the feed solution at a concentration of \qty{10}{\gram\per\litre}, and to the reactor whenever the pH dropped lower than \num[round-mode = places, round-precision = 1]{6}.

In the \Cref{eq:HRT,eq:OLR}, about \glsauto{hrt} and \glsauto{olr}, respectively:

\begin{equation}
    \mathrm{HRT} = \frac{V}{Q} \quad \left[\unit{\day}\right]
    \label{eq:HRT}
\end{equation}
\begin{equation}
    \mathrm{OLR} = \frac{S_{\text{in}} \times Q}{V} \quad \left[\unit{\kilo\gram\VS\per\cubic\meter\per\day}\right]
    \label{eq:OLR}
\end{equation}
where $S_{\text{in}}$ is the influent \unit{\VS} concentration (\unit{\kilo\gram\per\cubic\metre}), $Q$ is the volumetric flow rate (\unit{\cubic\metre\per\day}), and $V$ is the reactor working volume (\unit{\cubic\metre}).

\subsection{Feeding Solution}
\label{subsec:feeding_solution}

The feed was designed to achieve a \glsauto{cn_ratio} ratio of 25:1, which is optimal for \glsauto{ad} \parencite{li2011}. Characterisation of the \glsauto{fw} and \glsauto{cgl} showed that both substrates contained nitrogen, making it difficult to balance the \glsauto{cn_ratio} ratio using both of them together. Therefore, \glsauto{pgl} was used as an additive to allow precise control of the \glsauto{cn_ratio}. Based on the chosen \glsauto{hrt} and the daily volumetric flow rate, the reactor was fed once per day to ensure a flow rate high enough to transport the solid \glsauto{fw} suspension through the feeding tube.

The feed was made with a restriction of needing to have less than \qty{15}{\percent} of TS\unit{\percent}. The amount of \glsauto{olr} and if \glsauto{khco3} was added, was decided by the pH and how was the \glsauto{ch4} production of the reactor. If it was too acidic, it was used a lower \glsauto{olr} with \glsauto{khco3}, if it was at the right amount it was keep the same amount, and if it was too basic it was added a higher \glsauto{olr} without \glsauto{khco3}. The feed storage always had some remainig feed, and the new one was added on top, diluting or concentrating the feed going to the reactor.

The feed was also supplemented with nutrients to improve the growth of microorganisms. \Cref{tab:nutri} shows the nutrients added to the substrate and their respective concentrations.

\begingroup
\tablesetup
\footnotesize
\setlength{\tabcolsep}{3pt}
\begin{ThreePartTable}
    \begin{xltabular}{\linewidth}{@{} C C @{}}

        \caption{Nutrients added to the substrate and there dosages. (Adapted from: \textcite{safaa2025})}
        \label{tab:nutri} \\

        \toprule
        \thead{Nutrients} &
        \thead{Concentration (\unit{\milli\gram\per\litre})} \\
        \midrule
        \endfirsthead

        \multicolumn{2}{l}{\textit{\Cref{tab:nutri} continued}} \\
        \toprule
        \thead{Nutrients} &
        \thead{Concentration (\unit{\milli\gram\per\litre})} \\
        \midrule
        \endhead

        \bottomrule
        \endlastfoot

        \csvreader[
            late after line=\\\addlinespace
        ]{data/processed/nutrients.csv}{}
        {\expandafter\ce\expandafter{\csvcoli} & \num{\csvcolii}}

    \end{xltabular}
\end{ThreePartTable}
\endgroup

\subsection{Biogas Production}
\label{subsec:biogas_production}

The \glsauto{ampts} system (hardware~\Cref{fig:ampts-hard}, software~\Cref{fig:ampts-soft}), was provided to measure the gas flow produced by the reactor, with a \glsauto{co2} fixing unit using a solution of \qty{3}{\molar\NaOH}.

\begin{figure}[H]
    \centering
    \begin{subfigure}[t]{0.45\textwidth}
        \centering
        \includegraphics[width=\linewidth]{ampts.jpg}
        \caption{\Glsentrylongauto{ampts} gas flow messurer}
        \label{fig:ampts-col}
    \end{subfigure}%
    \hspace{1em}
    \begin{subfigure}[t]{0.45\textwidth}
        \centering
        \includegraphics[width=\linewidth, height=0.35\textheight, keepaspectratio]{ampts_co2.jpg}
        \caption{\Glsentrylongauto{ampts} \glsentrysymbolauto{co2} trap}
        \label{fig:ampts-co2}
    \end{subfigure}
    \caption{\Glsentrylongauto{ampts} hardware}
    \label{fig:ampts-hard}
\end{figure}

\begin{figure}[H]
    \centering
    \includegraphics[width=0.8\linewidth]{ampts_sys-print.png}
    \caption{\Glsentrylongauto{ampts} software}
    \label{fig:ampts-soft}
\end{figure}

\subsection{Analysis Methods}
\label{subsec:analysis_methods}

The method of DiLallo \& Albertson \apud{dilallo1961}{ribas2007}, to measure the \glsauto{alk} and \glsauto{vfa}. The sample collected was first centrifuged, at \qty{2604}{\gforce} for \num{30} minutes, to filter the particulates before doing the measurements.

The volumes were measured in \unit{\milli\litre}, the sample volume was always \qty{50}{\milli\litre}, the normality (\unit{\normality}) of the acid and base, at first was \num{0.1}, but then it was changed to \num[round-mode = places, round-precision = 1]{1}. The \Cref{eq:TA} is expressed in \unit{\milli\gram\CaCO\per\litre}, and the \Cref{eq:VFA} is expressed in \unit{\milli\gram\CHCOOH\per\litre}.

\begin{equation}
    \mathrm{TA} = \frac{N_\text{acid} * V_\text{acid} * 50000}{V_\text{sample}} \quad \left[\unit{\milli\gram\CaCO\per\litre}\right]
    \label{eq:TA}
\end{equation}
\begin{equation}
    \mathrm{VFA} = \frac{N_\text{base} * V_\text{base} * 60000}{V_\text{sample}} \quad \left[\unit{\milli\gram\CHCOOH\per\litre}\right]
    \label{eq:VFA}
\end{equation}

\subsection{Feeding of the reactor}
\label{subsec:feed_reac}

\Cref{tab:feeds} shows all substrates used to feed the reactor and when they were made.

\begingroup
\sisetup{}
\footnotesize
\setlength{\tabcolsep}{3pt}
\begin{ThreePartTable}
    \begin{xltabular}{\linewidth}{@{} C C C C C C C @{}}

    \caption{Feed used throughout the experimental period}
    \label{tab:feeds} \\

    \toprule
    \thead{Day} &
    \thead{Substrate 1} &
    \thead{Subestrate 2} &
    \thead{Substrate 3} &
    \thead{Extra\\ Addition} &
    \thead{\GLS{olr}\\ \& \GLS{ts}\%} &
    \thead{Volume\\ Made} \\
    \midrule
    \endfirsthead

    \multicolumn{6}{l}{\textit{\Cref{tab:feeds} continued}} \\
    \toprule
    \thead{Day} &
    \thead{Substrate 1} &
    \thead{Subestrate 2} &
    \thead{Substrate 3} &
    \thead{Extra\\ Addition} &
    \thead{\GLS{olr}\\ \& \GLS{ts}\%} &
    \thead{Volume\\ Made} \\
    \midrule
    \endhead

    \bottomrule
    \endlastfoot

    \num{0}\tnote{a} &
    \glsauto{fw}, \qty{156}{\gram} &
    \glsauto{glucose}, \qty{1.9}{\gram} &
    &
    \glsauto{ca_hco3}, \qty{5}{\gram} &
    \num{1.5} \& \num{15} &
    \qty{500}{\milli\litre} \\
    \addlinespace
    \num{4} &
    \glsauto{cgl}, \qty{23.2}{\milli\litre} &
    \glsauto{pgl}\tnote{b}, \qty{31.5}{\milli\litre} &
    &
    \glsauto{khco3}, \qty{1.9}{\gram} &
    \num{0.8} \& \num{10} &
    \qty{250}{\milli\litre} \\
    \addlinespace
    \num{11} &
    \glsauto{fw}, \qty{12.54}{\gram} &
    \glsauto{cgl}, \qty{9.02}{\milli\litre} &
    \glsauto{pgl}, \qty{2.6}{\milli\litre} &
    \glsauto{khco3}, \qty{1.57}{\gram} &
    \numrange{0.13}{0.4} \& \numrange{5}{10}\tnote{c} &
    \qty{250}{\milli\litre} \\
    \addlinespace
    \num{18} &
    \glsauto{fw}, \qty{10.88}{\gram} &
    \glsauto{cgl}, \qty{2.89}{\milli\litre} &
    \glsauto{pgl}, \qty{1.44}{\milli\litre} &
    &
    \num{0.0256} \& \num{2} &
    \qty{250}{\milli\litre} \\
    \addlinespace
    \num{25} &
    \glsauto{fw}, \qty{31.12}{\gram} &
    \glsauto{cgl}, \qty{44}{\milli\litre} &
    \glsauto{pgl}, \qty{29}{\milli\litre} &
    &
    \num{1.4} \& \num{10} &
    \qty{250}{\milli\litre} \\
    \addlinespace
    \num{32} &
    \glsauto{fw}, \qty{31.01}{\gram} &
    \glsauto{cgl}, \qty{44}{\milli\litre} &
    \glsauto{pgl}, \qty{28.9}{\milli\litre} &
    &
    \num{1.4} \& \num{10} &
    \qty{250}{\milli\litre} \\
    \addlinespace
    \num{39} &
    \glsauto{fw}, \qty{46.52}{\gram} &
    \glsauto{cgl}, \qty{66}{\milli\litre} &
    \glsauto{pgl}, \qty{43}{\milli\litre} &
    \glsauto{khco3}, \qty{2.6}{\gram} &
    \num{3.1} \& \num{15} &
    \qty{250}{\milli\litre} \\
    \addlinespace
    \num{46} &
    \glsauto{fw}, \qty{50.07}{\gram} &
    \glsauto{cgl}, \qty{66}{\milli\litre} &
    \glsauto{pgl}, \qty{43}{\milli\litre} &
    \glsauto{khco3}, \qty{2.67}{\gram} &
    \num{3.16} \& \num{15} &
    \qty{250}{\milli\litre} \\
    \addlinespace
    \num{53}\tnote{d} &
    \glsauto{cgl}, \qty{55}{\milli\litre} &
    \glsauto{pgl}, \qty{11.5}{\milli\litre} &
    &
    \glsauto{khco3}, \qty{5.38}{\gram} &
    \num{0.42} \& \num{10} &
    \qty{500}{\milli\litre} \\
    \addlinespace
    \num{60} &
    \glsauto{cgl}, \qty{55}{\milli\litre} &
    \glsauto{pgl}, \qty{12}{\milli\litre} &
    &
    \glsauto{khco3}, \qty{5.18}{\gram} &
    \num{0.42} \& \num{10} &
    \qty{500}{\milli\litre}\tnote{e} \\
    \addlinespace
    \num{67} &
    \glsauto{cgl}, \qty{41}{\milli\litre} &
    \glsauto{pgl}, \qty{9}{\milli\litre} &
    &
    \glsauto{khco3}, \qty{2.56}{\gram} &
    \num{0.95} \& \num{15} &
    \qty{250}{\milli\litre} \\
    \addlinespace
    \num{74} &
    \glsauto{cgl}, \qty{60}{\milli\litre} &
    \glsauto{pgl}, \qty{12.5}{\milli\litre} &
    &
    \glsauto{khco3}, \qty{2.57}{\gram} &
    \num{1.37} \& \num{15} &
    \qty{250}{\milli\litre} \\

    \end{xltabular}

    \begin{tablenotes}
        \item[a] It only started feeding on day \num{2}.
        \item[b] The solution of \Gls{pgl} was diluted \qty{100}{\dil}.
        \item[c] The \glsauto{fw} wasn't well mixtured this day.
        \item[d] Change of substrate.
        \item[e] There was a leak on the feeding tube.
    \end{tablenotes}
\end{ThreePartTable}
\endgroup
